% =============================================================================
%  research-paper.tex — borrador de paper academico con la clase
%  claudeeditorial. Ejercita: ClaudePaperCover, abstract, keywords, MSC,
%  PaperAuthor (multi-autor con afiliacion), theorems en castellano e
%  ingles compartiendo contador, demostraciones, \cref, \norm, \abs, \E,
%  \rank, \argmin, \argmax, eqs numeradas en terracota, conjeturas y
%  observaciones.
%
%  Compilacion:
%    xelatex examples/research-paper.tex   (dos pasadas para \cref)
%
%  El contenido es ficticio; se usa para mostrar el sistema visual del
%  template aplicado a un working paper.
% =============================================================================
\documentclass{claudeeditorial}

\renewcommand{\DocKicker}{Working paper · v1}
\renewcommand{\DocTitle}{Coeficiente de regresión array (AR): formalización rigurosa y delimitación de las conjeturas abiertas}
\renewcommand{\DocSubtitle}{Notas y borrador del estudio para un journal aplicado: definiciones canónicas, propiedades demostradas y casos pendientes}
\renewcommand{\DocAuthor}{Equipo AR · Lorem Ipsum}
\renewcommand{\DocRole}{Working paper}
\renewcommand{\DocDate}{5 de mayo de 2026}
\renewcommand{\DocRef}{WP-AR-001 · vol. I}
\renewcommand{\DocFooter}{claudeeditorial · research paper}

% Metadatos del paper
\renewcommand{\PaperKeywords}{coeficiente AR; estimadores robustos; calibración exacta; rankings; permutaciones; footrule}
\renewcommand{\PaperMSC}{60E15, 62G35, 65C60}

\PaperAuthor{Lorem Ipsum}{Departamento de Matemática Aplicada, Universidad de Lorem}{lorem@uni.example}
\PaperAuthor{Dolor Sit Amet}{Centro de Investigación Editorial, Universidad de Dolor}{dolor@uni.example}
\PaperAuthor{Consectetur Adipiscing}{Departamento de Estadística, Universidad de Consectetur}{ca@uni.example}

\hypersetup{
  pdftitle={Coeficiente AR: formalizacion rigurosa},
  pdfauthor={Lorem Ipsum, Dolor Sit Amet, Consectetur Adipiscing},
  pdfsubject={Working paper - claudeeditorial}
}

% Macros propias del paper (la clase ya provee \E, \R, \N, \rank, ...)
\newcommand{\ARop}{\operatorname{AR}_{[0,1]}}
\newcommand{\ARc}{\operatorname{AR}_{[-1,1]}}
\newcommand{\Perm}{\Pi}

\ClaudeRunningHeader
\setstretch{1.20}

\begin{document}
\BodyCopy

\ClaudePaperCover{%
  Se formaliza el coeficiente AR (\emph{Array Regression Coefficient})
  y se demuestran con rigor sus propiedades operativas: la identidad
  probabilística exacta $D(r,n)=\E\abs{r-S}$ para $S$ uniforme en
  $\set{0,\ldots,n-1}$, la calibración $\E[\ARop]=\tfrac{1}{2}$ bajo
  permutación uniforme \emph{para todo $n\ge 2$, sin pasar al límite
  asintótico}, y la suma cerrada $\sum_{r=0}^{n-1}D(r,n)=(n^2-1)/3$.
  Se separa explícitamente lo demostrado de lo conjeturado: la
  afirmación de que el mínimo de $\ARop$ sobre $\Perm(n)$ se alcanza
  en la rotación cíclica de paso $\lfloor n/2\rfloor$ queda como
  conjetura, verificada por enumeración exhaustiva para $n\le 8$.
  Toda constante numérica del cuerpo principal proviene de aritmética
  racional exacta.}

\tableofcontents
\clearpage

\ClaudeChapterPage{01}{Marco y notación}{Convenciones temporales, espacio de permutaciones y el denominador ordinal $D(r,n)$ que articula toda la construcción.}

\section{Notación}\label{sec:notacion}

Sea $E=\set{1,\ldots,n}$ un conjunto finito de elementos comparables y
$t=1,\ldots,T$ una discretización del tiempo. Se observa un valor
$x_{i,t}\in\R$ para cada elemento $i\in E$ y periodo $t$. Para
$z\in\R^n$, $\rank(z)_i\in\set{0,\ldots,n-1}$ es el rango ascendente
de $z_i$ con desempate por índice.

\begin{definition}[denominador ordinal]\label{def:D}
Para $r\in\set{0,\ldots,n-1}$ definimos
\begin{equation}\label{eq:D}
  D(r,n) := \frac{r^2+r}{n}+\frac{n}{2}-r-\frac{1}{2}.
\end{equation}
\end{definition}

\begin{proposition}[positividad y simetría]\label{prop:D-basico}
Para todo $n\ge 2$ y $r\in\set{0,\ldots,n-1}$ se cumple
$D(r,n)>0$ y $D(r,n)=D(n-1-r,n)$. Además
\begin{equation}\label{eq:sumD}
  \sum_{r=0}^{n-1}D(r,n)=\frac{n^2-1}{3}.
\end{equation}
\end{proposition}

\begin{proof}
La simetría se verifica sustituyendo $n-1-r$ en \cref{eq:D}. Tratando
$r$ como variable real, $D'(r,n)=(2r+1)/n-1$, que se anula en
$r^\ast=(n-1)/2$. Allí $D(r^\ast,n)=(n^2-1)/(4n)>0$. Como $D$ es
convexa, ese es el mínimo real. Para \cref{eq:sumD},
\[
  \sum_{r=0}^{n-1}D(r,n)
  =\frac{1}{n}\sum_{r=0}^{n-1}(r^2+r)
   +\sum_{r=0}^{n-1}\!\left(\tfrac{n}{2}-r-\tfrac{1}{2}\right).
\]
El segundo sumatorio se anula y el primero da $(n^2-1)/3$.
\end{proof}

\begin{proposition}[interpretación probabilística]\label{thm:D-prob}
Si $S\sim\mathcal{U}\set{0,\ldots,n-1}$, entonces
$D(r,n)=\E\abs{r-S}$.
\end{proposition}

\begin{proof}
Por definición,
$\E\abs{r-S}=\frac{1}{n}\sum_{s=0}^{n-1}\abs{r-s}
=\frac{1}{n}\bigl(\sum_{k=0}^{r}k+\sum_{k=1}^{n-1-r}k\bigr)$.
Usando $\sum_{k=0}^{m}k=m(m+1)/2$ se obtiene el lado derecho de
\cref{eq:D}.
\end{proof}

\begin{observacion}
\Cref{thm:D-prob} permite leer $D(r,n)$ como la \emph{distancia
esperada} de un rango fijo $r$ a un rango uniforme. Esa interpretación
es la base operativa del normalizador en \cref{eq:ar-op}.
\end{observacion}

\section{Definición del coeficiente}

Sean $p,q$ permutaciones de $\set{0,\ldots,n-1}$. Pensamos $p$ como
predicción y $q$ como observación.

\begin{definition}[escalas operativa y centrada]\label{def:dos-escalas}
\begin{align}
  \ARop(p,q) &:= 1-\frac{1}{2n}\sum_{i=1}^{n}\frac{\abs{q_i-p_i}}{D(q_i,n)}, \label{eq:ar-op}\\
  \ARc(p,q)  &:= 2\,\ARop(p,q)-1.\label{eq:ar-centered}
\end{align}
\end{definition}

\begin{remark}
Bajo azar uniforme, $\E[\ARop]=\tfrac{1}{2}$ y $\E[\ARc]=0$. La escala
operativa es cómoda para tablas porque pone el azar en $0.5$; la
escala centrada es la que aparece en el código de referencia.
\end{remark}

\begin{theorem}[calibración exacta]\label{thm:calibracion}
Sea $P\sim\mathcal{U}(\Perm(n))$ y $R\in\Perm(n)$ fijo. Entonces
$\E[\ARop(R,P)]=\tfrac{1}{2}$ \emph{para todo $n\ge 2$}.
\end{theorem}

\begin{proof}
Por linearidad del valor esperado y \cref{thm:D-prob},
\[
  \E[\ARop(R,P)]
  =1-\frac{1}{2n}\sum_{i=1}^{n}\frac{\E\abs{P_i-R_i}}{D(R_i,n)}
  =1-\frac{1}{2n}\sum_{i=1}^{n}1
  =\frac{1}{2}. \qedhere
\]
\end{proof}

\begin{conjetura}[mínimo de AR]\label{conj:argmin}
Para todo $n\ge 3$, el mínimo de $\ARop(\mathrm{id},\sigma)$ sobre
$\sigma\in\Perm(n)$ se alcanza en la rotación cíclica de paso
$\lfloor n/2\rfloor$.
\end{conjetura}

\begin{nota}
\Cref{conj:argmin} se ha verificado por enumeración exhaustiva para
$n\le 8$, sin demostración combinatoria general. La rotación de paso
$\lfloor n/2\rfloor$ tiene límite asintótico $\ARop\to 1-\pi/4$.
\end{nota}

\section{Indicadores y trazabilidad}

Para mantener trazable el cuerpo principal del paper, registramos los
indicadores reproducibles del estudio.

\begin{tcbraster}[raster columns=4,raster equal height=rows,raster column skip=8pt]
\KPICard{Teoremas}{4}{con demostración completa}
\KPICard{Conjeturas}{1}{verificada para $n\le 8$}
\KPICard{Verificación}{Exacta}{aritmética racional}
\KPICard{Réplica}{Pendiente}{cohorte $n=12$}
\end{tcbraster}

\section{Decisión metodológica}

\ClaudeDecisionRecord[Activa]
  {ADR-001 · Aritmética racional para constantes del cuerpo}
  {Las constantes numéricas que aparecen en teoremas y observaciones provienen, históricamente, de cálculos en coma flotante con redondeo no controlado.}
  {Toda constante del cuerpo se computa en aritmética racional exacta (módulo \texttt{fractions}). Los resultados decimales aparecen sólo en el apéndice numérico.}
  {Coste: una pasada de verificación adicional. Beneficio: ningún valor del cuerpo proviene de redondeo, lo que neutraliza una clase entera de objeciones.}

\section{Hipótesis abiertas}

\begin{ClaudeTable}{@{}lL L l@{}}
\ClaudeTableHeader Hipótesis & Predicción & Validación pendiente & Estado \\
\toprule
\ClaudeRiskRow{H-001}{El mínimo de $\ARop$ se alcanza en la rotación de paso $\lfloor n/2\rfloor$ para todo $n\ge 3$.}{Demostración combinatoria general (enumeración hasta $n=8$).}{\StatusOpen}
\ClaudeRiskRow{H-002}{$\norm{\ARc(p,q)-\ARc(q,p)}_\infty = O(1/n)$.}{Cota empírica vía simulación; cota formal abierta.}{\StatusWIP}
\bottomrule
\end{ClaudeTable}

\section{Verificación numérica}

La verificación se hizo en aritmética racional exacta, sin punto
flotante en el cuerpo principal.

\begin{ClaudeCode}[title={Verificación racional exacta del coeficiente AR}]{Python}
from fractions import Fraction

def D(r, n):
    return Fraction(r*r + r, n) + Fraction(n, 2) - r - Fraction(1, 2)

def AR_op(p, q):
    n = len(p)
    s = sum(Fraction(abs(q[i] - p[i])) / D(q[i], n) for i in range(n))
    return 1 - s / (2 * n)

p = list(range(8))             # identidad
q = [4, 5, 6, 7, 0, 1, 2, 3]   # rotacion de paso n/2
print(AR_op(p, q))             # Fraction(67, 280)
\end{ClaudeCode}

\begin{ClaudeOutput}
\texttt{Fraction(67, 280)} — equivalente a \(0.2392857\ldots\); no
proviene de coma flotante. \Cref{conj:argmin} predice este valor
como mínimo sobre $\Perm(8)$.
\end{ClaudeOutput}

\section{Conclusiones}

Hemos demostrado las propiedades operativas de $\ARop$
(\cref{thm:D-prob,thm:calibracion,prop:D-basico}) y separado lo
\emph{demostrado} de lo \emph{conjeturado} (\cref{conj:argmin}). El
borrador queda listo para revisión interna; la siguiente iteración
ataca la demostración combinatoria del mínimo de AR.

\ClaudeSignature{\DocAuthor}{\DocRole}

\end{document}
